% !TEX root = ../main.tex

\section{Discusi\'on}
\label{sec:discussion}

Los resultados de la secci\'on~\ref{sec:results} dibujan una imagen ambivalente
que merece discutirse con cuidado. Por un lado, \sculpt es t\'ecnicamente
exitoso: el lenguaje es Turing-completo, su int\'erprete funciona, las
piezas f\'isicas resisten el uso y la mayor\'ia de los participantes
identifica y manipula los componentes sin dificultad. Por otro,
escribir programas en \sculpt sigue siendo dif\'icil. Aproximadamente
la mitad de los participantes completa las tareas computacionales,
y los programadores expertos no superan a los adolescentes novatos
en esa m\'etrica.

Esta tensi\'on no es un fracaso del lenguaje puntualmente, sino es una de sus
caracter\'isticas constitutivas. \sculpt no busca eliminar la
dificultad intr\'inseca de programar, ese problema lo comparten todos
los lenguajes. La intenci\'on cambiar la \emph{naturaleza} de esa dificultad y,
al hacerlo, ampliar qui\'enes la encuentran interesante. Lo que los
resultados muestran es que la dificultad se reasigna: cae con fuerza la
barrera de acceso a la primera interacci\'on con el lenguaje, los
ni\~nos manipulan piezas y los artistas las disfrutan. Se conserva,
casi intacta, la dificultad de razonar algor\'itmicamente sobre pilas, condicionales y bucles.
La discrepancia entre la percepci\'on t\'ecnica baja de los
participantes (medias Likert 2{,}2--3{,}4) y la percepci\'on docente
alta (4{,}0--4{,}3) refuerza esta lectura. \sculpt opera mejor como
veh\'iculo did\'actico mediado por un facilitador que como herramienta
de programaci\'on aut\'onoma.

La fisicalidad y estetica merece tratarse como objeto de estudio en s\'i misma, no
solo como medio. Las sesiones con artistas son el caso m\'as claro de esto.
Participantes que reportaron mayor satisfacci\'on est\'etica son
tambi\'en los que m\'as alejaron las piezas de su uso computacional
prescrito, construyendo esculturas no funcionales y explorando
ensambles no previstos en la especificaci\'on. El lenguaje, en sus
manos, dej\'o de ser una herramienta para volverse un material. Esta
observaci\'on conecta directamente con preguntas abiertas en la
comunidad de \emph{programming experience}: qu\'e se siente programar
en un medio que se puede sostener, qu\'e prop\'ositos creativos
emergen cuando el c\'odigo deja la pantalla.

\paragraph{Posicionamiento.}
\sculpt no aspira a reemplazar los lenguajes textuales, ni ser\'ia
prudente que lo intentara. Su aparatosa naturaleza para proyectos de software
serios es insantaneamente aparente. El tama\~no de un programa est\'a acotado por la
cantidad de piezas que f\'isicamente se pueden ensamblar, el costo de
impresi\'on crece linealmente con esa cantidad, el ensamblaje toma
tiempo y, al ser solo fisico, ni siquiera corre el resultado ensamblado en un computador.
\sculpt funciona, en cambio, como puerta
de entrada para audiencias hist\'oricamente alejadas de la
programaci\'on y como medio de expresi\'on para artistas que quieren
incorporar c\'omputo en su pr\'actica sin pasar por el teclado.

\paragraph{Limitaciones.}
Adem\'as de las restricciones de tama\~no ya mencionadas, conviene
nombrar tres l\'imites m\'as. La primera y m\'as grande obedece a la naturaleza f\'isica del ejercicio.
Las sesiones de validaci\'on requieren de facilitadores que gu\'ian a los participantes,
resuelven dudas, arman esculturas, supervisan el uso de las piezas y, en el caso de los menores de edad,
se aseguran de que los docentes acompa\~nen a los estudiantes. Esto hace que la escalabilidad
del ejercicio sea limitada, y que su aplicaci\'on en contextos educativos dependa de la 
disponibilidad de facilitadores capacitados.
La segunda limitaci\'on va de la mano con la primera. El uso de piezas impresas en 3D con acoples
magn\'eticos es una soluci\'on pr\'actica para prototipar y manufacturar el lenguaje a un costo accesible. Sin embargo,
el material es fr\'agil y se desgasta con el uso, reduciendo el inventario de cada sesi\'on.
Adicionalmente, el costo temporal de imprimir piezas suficientes para dictar los talleres es elevado con una 
disponibilidad limitada de impresoras 3D.
Esto limita mucho la cantidad de participantes que se pueden atender en una sesion.
Esto es evidente en los resultados: el n\'umero de participantes es bajo, y el n\'umero de tareas computacionales asignadas es limitado.
Esta limitaci\'on se puede superar con materiales m\'as resistentes, procesos de impresi\'on
masivos y optimización del dise\~no.
Tercero, las muestras de validaci\'on, aunque diversas, no son grandes: el desempe\~no
educativo a m\'as largo plazo, en cohortes estables, sigue siendo
una pregunta abierta.

\endinput
